%TODO
\section{牛顿法}

牛顿法实质上是把非线性方程\(f(x) = 0\)近似地转化为一个线性方程.

简化牛顿法：
基于假设“小邻域内导数变化不大”
用\(f'(x_0)\)代替\(f'(x_k)\)

牛顿下山法：
当\(k\)充分大时，
\(f(x_{k+1})\)与\(f(x_k)\)的差别不大，
\(f'(x_{k+1})\)与\(f'(x_k)\)的差别不大，
这时收敛速度较慢.
为了加速收敛，
使用一个因子\(\lambda\)


对于\(m\)重根，
\(f(x) = (x - x^*)^m g(x)\)，
牛顿法不是平方收敛的，
可将迭代法改为
\begin{equation*}
	x_{k+1} = x_k - m \frac{f(x_k)}{f'(x_k)}.
\end{equation*}

